\section{Conclusions and Related Work}\label{sec:conclusion}

We have introduced MEC, a polarization measure defined as the minimum
effort required to bring an opinion distribution to consensus. The
measure rests on three connected mathematical foundations: a
transportation-theoretic interpretation as a Wasserstein-type distance
to the set of consensus distributions, a Minkowski-type characterization
as a weighted $L^\beta$-dispersion of the $\alpha$-power distribution,
and an axiomatic compatibility with the framework of Esteban and
Ray~\cite{esteban1994measurement}. The basic case $(\alpha,\beta)=(1,1)$
recovers the Earth Mover's Distance to the extremal distribution, and
the cases $(1,1)$ and $(1,2)$ specialize to mean absolute deviation
about the median and to a variance-like dispersion respectively. The
parameters $\alpha$ and $\beta$ admit a natural reading in terms of
group identification~\cite{TajfelTurner1979} and cognitive cost of
belief revision~\cite{Festinger1957}.

Beyond the scalar value, the optimization formulation produces an
optimal consensus point $y^*$, which existing pairwise and dispersion
measures do not provide. This additional structure has both
theoretical and practical use. On the theoretical side, $y^*$ anchors
the proofs of the Minority Principle and of the shifting-away
result, and supports a closed-form Tipping Point method that locates
the critical mass at which a partial transfer of opinions away from
$y^*$ ceases to reduce polarization and begins to increase it. The
analysis also reveals that polarization is not monotone in extremism,
nor in movement away from the optimal consensus point: whether such a
movement raises or lowers polarization depends on whether a whole group
or only a fragment of it is displaced. Empirically, MEC matches the
strongest parametrization of \ERLabel{} on the expert-judgment
benchmark of Koudenburg, Kiers, and Kashima~\cite{koudenburg2021new}, and
outperforms the dispersion-based alternatives \VdELabel{} and
\TWDLabel.

The proposal is part of a broader effort to bring tools from optimal
transport and convex analysis to the quantitative study of polarization,
a field that has been dominated by pairwise constructions in the spirit
of Gini and Esteban--Ray~\cite{esteban1994measurement,foster85}. The
Earth Mover's Distance has appeared in this literature as a measure of
distance between distributions~\cite{geher2020isometric,vallender1974calculation},
but, to our knowledge, has not previously been used to define
polarization itself as a distance to a structured target set. The
parametric framework $(\alpha,\beta)$ also relates to the
inequality-measurement tradition~\cite{foster85}, in which similar
exponents control sensitivity to mass and distance, and to the
sociological literature on group identification and cognitive
dissonance~\cite{TajfelTurner1979,Festinger1957}, which motivates the
specific role of the two parameters.

Several directions remain open. The current formulation treats opinions
as elements of $[0,1]$ and distributions as having finite support. A
natural extension is to general metric spaces, in which the ground
distance $|x-y|$ would be replaced by an arbitrary metric and the
domain by a suitable measurable space. The parametric family invites
further study of the dependence of $\MEC$ on $(\alpha,\beta)$, in
particular a theoretical justification of the empirically successful
parametrization $(2,1.15)$. The Tipping Point method, established here
for $\alpha=\beta=2$, may admit extensions to other parameter ranges.
The empirical validation rests on a synthetic benchmark; testing MEC
on real survey data, and on dynamic models of opinion update, would
clarify both its descriptive accuracy and its sensitivity to noise.
Finally, the optimization formulation suggests a connection to
mediation and policy framing through the optimal consensus point
$y^*$, which we have treated here only as a structural object but
which may admit substantive interpretation in applied settings.

\paragraph{Acknowledgments.}
The authors would like to thank the anonymous reviewers for their very
insightful comments on an earlier version of this paper.  This work
has been partially supported by the SGR project PROMUEVA (BPIN
2021000100160) under the supervision of Minciencias (Ministerio de
Ciencia Tecnolog\'ia e Innovaci\'on, Colombia),  by the CNRS project TOBIAS under the MITI interdisciplinary program, and the ECOSNORD Project BRAINS C26M02. 
